\documentclass[../../main.tex]{subfiles}
\begin{document}
% SOURCE p=223; printed=208; object=ch06-interpolation-exercises-title; region=upper-middle
\subsection*{习题}
\setcounter{exer}{40}
% SOURCE p=223; printed=208; object=ch06-ex41; number=41; region=middle
\begin{exercise}\label{exercise:ch06-symmetric-operator-interpolation}
设 $1<p\leqslant\infty$、$p^{-1}+q^{-1}=1$。若 $T$ 为 $L^p(\mathbb R^n)$ 上的有界算子，并对所有 $f,g\in L^p\cap L^q$ 满足 $\int(Tf)g=\int f(Tg)$，则 $T$ 唯一延拓为每个 $L^r$ 上的有界算子，其中当 $p<q$ 时 $r\in[p,q]$，当 $q<p$ 时 $r\in[q,p]$。
\end{exercise}
% SUPPLEMENT: id=v14-ch06-ex41-solution; parent=ch06-ex41; kind=solution+motivation+hypotheses; origin=independent-derivation
\begin{solution}
\textbf{由配对对称性建立另一端点（重构新增）。}这里有界算子按本章习惯指有界线性算子。记 $M=\|T\|_{L^p\to L^p}$、$D=L^p\cap L^q$。对 $f\in D$ 与有限测度支集的简单函数 $g$，有 $g\in D$，由原假设及 Hölder 不等式，
\begin{align*}
\left|\int(Tf)g\right|=\left|\int f(Tg)\right|
\leqslant\|f\|_q\|Tg\|_p
\leqslant M\|f\|_q\|g\|_p.
\end{align*}
左边的乘积可积，因为 $Tf\in L^p$ 且 $g\in L^q$。$\mathbb R^n$ 的 Lebesgue 测度为 $\sigma$ 有限，由\cref{theorem:ch06-dual-test-criterion}，$Tf\in L^q$ 且 $\|Tf\|_q\leqslant M\|f\|_q$。由于 $p>1$，共轭指数 $q<\infty$，有限测度支集的简单函数包含于 $D$ 并在 $L^q$ 稠密，故 $T|_D$ 唯一延拓为有界线性映射 $S:L^q\to L^q$，范数不超过 $M$，且 $S=T$ 于 $D$。

在和空间上定义 $U(a+b)=Ta+Sb$，$a\in L^p,b\in L^q$。若另有 $a+b=a'+b'$，则 $a-a'=-(b-b')\in D$，所以 $T(a-a')=S(a-a')=-S(b-b')$；故 $U$ 定义良好。它在两个端点均有范数界 $M$。由\cref{theorem:ch06-riesz-thorin}，每个严格中间指数 $r$ 都有 $\|Uf\|_r\leqslant M\|f\|_r$，端点分别为原来的 $T$ 与 $S$。若 $p=q=2$，没有中间区间，直接使用原算子即可。

最后证明唯一性。对于有限 $r$，有限测度支集的简单函数在 $L^r$ 稠密，且位于 $D$，任何与原 $T$ 相容的有界延拓都必须在这些简单函数上相同，因而在整个 $L^r$ 相同。可能的无穷端点只能是原输入端点 $p=\infty$，那里算子已经给定。于是全部所述延拓唯一。
\end{solution}
\begin{remark}
\textbf{动机与逐条件说明（重构新增）。}配对对称性把对 $Tg$ 的 $L^p$ 控制转移成对 $Tf$ 的 $L^q$ 控制，然后由插值填满两个指数之间的区间。配对按原题为双线性积分，没有额外复共轭；Hölder 只估计绝对值，所以并不要求 $T$ 是某个 Hilbert 空间上的自伴算子。$p>1$ 保证新端点 $q$ 有限，以便由简单函数稠密性获得唯一延拓；$p=\infty,q=1$ 也完全包含在证明内。$\mathbb R^n$ 的 $\sigma$ 有限测度保证简单测试恢复对偶范数。无限测度下各 $L^r$ 一般互不包含，题目中的“延拓”应理解为在交集上一致、由同一个和空间算子诱导；不能理解成原定义域必是每个新空间的子集。
\end{remark}
% END SUPPLEMENT: id=v14-ch06-ex41-solution
% SOURCE p=223; printed=208; object=ch06-ex42; number=42; region=middle
\begin{exercise}\label{exercise:ch06-marcinkiewicz-equal-input}
证明 Marcinkiewicz 定理在 $p_0=p_1$ 时的情形。（令 $p=p_0=p_1$，有 $\lambda_{Tf}(\alpha)\leqslant(C_0\|f\|_p/\alpha)^{q_0}$，也有 $\lambda_{Tf}(\alpha)\leqslant(C_1\|f\|_p/\alpha)^{q_1}$。根据 $\alpha$ 选择较好的一个界，用来控制 $q\int_0^\infty\alpha^{q-1}\lambda_{Tf}(\alpha)\,\mathrm d\alpha$。）
\end{exercise}
% SUPPLEMENT: id=v14-ch06-ex42-solution; parent=ch06-ex42; kind=solution+motivation+hypotheses; origin=independent-derivation
\begin{solution}
\textbf{比较两个尾部界（重构新增）。}由 $q_0\ne q_1$、$p\leqslant q_j$，$p$ 必有限。交换端点编号后可设 $q_0<q_1$，所以 $q_0<q<q_1$，且 $q<\infty$。零输入由端点估计直接处理；非零输入借正齐次性归一为 $\|f\|_p=1$。先设 $q_1<\infty$，并选正的端点常数。两条界相交于
\begin{align*}
a=\left(\frac{C_1^{q_1}}{C_0^{q_0}}\right)^{1/(q_1-q_0)}.
\end{align*}
当 $0<\alpha<a$ 时使用 $C_0^{q_0}\alpha^{-q_0}$，当 $\alpha>a$ 时使用 $C_1^{q_1}\alpha^{-q_1}$。由层积分公式，
\begin{align*}
\|Tf\|_q^q
&\leqslant qC_0^{q_0}\int_0^a\alpha^{q-q_0-1}\,\mathrm d\alpha
+qC_1^{q_1}\int_a^\infty\alpha^{q-q_1-1}\,\mathrm d\alpha\\
&=q\left[\frac{C_0^{q_0}a^{q-q_0}}{q-q_0}
+\frac{C_1^{q_1}a^{q-q_1}}{q_1-q}\right]<\infty.
\end{align*}
其 $1/q$ 次方可取为 $B_{p,q}$，缩放恢复 $\|Tf\|_q\leqslant B_{p,q}\|f\|_p$。

若 $q_1=\infty$，弱端点约定给出 $\|Tf\|_\infty\leqslant C_1$，所以 $\lambda_{Tf}(\alpha)=0$ 对 $\alpha\geqslant C_1$ 成立。于是
\begin{align*}
\|Tf\|_q^q\leqslant qC_0^{q_0}\int_0^{C_1}\alpha^{q-q_0-1}\,\mathrm d\alpha
=\frac{qC_0^{q_0}C_1^{q-q_0}}{q-q_0}<\infty.
\end{align*}
再取根并缩放，即完成全部情形。
\end{solution}
\begin{remark}
\textbf{动机、参数依赖与逐条件说明（重构新增）。}相同输入指数时无需分解输入：同一个 $Tf$ 已同时拥有两种尾界，低高度选较小幂、高高度选较大幂即可。严格的中间指数保证 $q-q_0>0$、$q_1-q>0$，正是两个积分收敛的条件；无穷端点则以有界性截去高尾。这里没有交换原测度空间上的积分，因而无需 $\sigma$ 有限性；也没有实际使用次可加性，只使用两个端点界与缩放。

此情形中 $p$ 固定，常数一般仍依赖 $q$，所以用 $B_{p,q}$ 比原文的 $B_p$ 更明确。不能删除这种依赖：取输入空间为测度一的单点、$p=1$，输出空间为 $(1,\infty)$，并令 $Tc(x)=c/x$。它从 $L^1$ 到弱 $L^1$ 的界为一，且到 $L^\infty$ 的界也为一，但 $\|T1\|_q^q=1/(q-1)$ 对 $q>1$ 成立，故 $q\downarrow1$ 时范数无界。它满足这道题的两个端点条件，却不可能有与所有中间 $q$ 无关的统一常数。
\end{remark}
% END SUPPLEMENT: id=v14-ch06-ex42-solution
% SOURCE p=223; printed=208; object=ch06-ex43; number=43; region=middle
\begin{exercise}\label{exercise:ch06-maximal-indicator}
设 $H$ 为 $\mathbb R$ 上的 Hardy–Littlewood 极大算子。显式计算 $H\chi_{(0,1)}$；证明它对所有 $p>1$ 都属于 $L^p$，且属于弱 $L^1$ 而不属于 $L^1$；当 $p\to1$ 时，其 $L^p$ 范数像 $(p-1)^{-1}$ 一样趋于无穷，尽管 $\|\chi_{(0,1)}\|_p=1$ 对所有 $p$ 都成立。
\end{exercise}
% SUPPLEMENT: id=v14-ch06-ex43-solution; parent=ch06-ex43; kind=solution+motivation+hypotheses; origin=independent-derivation
\begin{solution}
\textbf{球平均的精确计算（重构新增）。}令 $f=\chi_{(0,1)}$、$h=Hf$，此处球为以 $x$ 为中心的区间 $(x-r,x+r)$。若 $0<x<1$，取足够小的 $r$ 可使平均为一，而任意平均都不超过一，故 $h(x)=1$。若 $x\leqslant0$，当 $r\leqslant-x$ 时平均为零；当 $-x<r\leqslant1-x$ 时平均为 $(x+r)/(2r)=1/2+x/(2r)$，它关于 $r$ 非减；当 $r\geqslant1-x$ 时平均为 $1/(2r)$，非增。最大值在 $r=1-x$ 取得。由关于 $1/2$ 的对称性，
\begin{align*}
h(x)=\begin{cases}
[2(1-x)]^{-1},&x\leqslant0,\\
1,&0<x<1,\\
(2x)^{-1},&x\geqslant1.
\end{cases}
\end{align*}
特别地，两端点的值都是 $1/2$，不是一。

对 $1<p<\infty$，
\begin{align*}
\|h\|_p^p=1+2\int_1^\infty(2x)^{-p}\,\mathrm dx
=1+\frac{2^{1-p}}{p-1}<\infty,
\end{align*}
而 $\|h\|_\infty=1$。当 $p=1$ 时尾积分为调和积分，故 $h\notin L^1$。根据上述分段式，其分布函数为
\begin{align*}
\lambda_h(\alpha)=\begin{cases}
\alpha^{-1}-1,&0<\alpha\leqslant1/2,\\
1,&1/2<\alpha<1,\\
0,&\alpha\geqslant1.
\end{cases}
\end{align*}
第一段的水平集两端为 $1-1/(2\alpha)$ 与 $1/(2\alpha)$，忽略有限个端点后长度为 $\alpha^{-1}-1$；其余两段由中间平台直接得到。因此
$[h]_1=\sup_{\alpha>0}\alpha\lambda_h(\alpha)=1$，故 $h$ 属于弱 $L^1$。

最后令 $\varepsilon=p-1\downarrow0$。由精确范数式，
\begin{align*}
\log\bigl(\varepsilon\|h\|_{1+\varepsilon}\bigr)
=\frac{\varepsilon\log\varepsilon+\log(\varepsilon+2^{-\varepsilon})}{1+\varepsilon}
\longrightarrow0.
\end{align*}
所以 $(p-1)\|h\|_p\to1$，这严格证明了题目的等价增长，而不只是证明 $p$ 次方范数的增长。输入 $f$ 的支集测度为一且非零值为一，故所有这些输入范数都等于一。
\end{solution}
\begin{remark}
\textbf{动机与逐条件说明（重构新增）。}一个有界紧支集输入，经过极大算子后产生两侧 $|x|^{-1}$ 的长尾；弱 $L^1$ 能容纳这个临界尾部，$L^1$ 积分却不能。使用的是有中心球的定义，若改成所有包含 $x$ 的区间，精确分段式和常数都会改变。维数一使球平均可以直接按长度求出；$p>1$ 恰是尾部积分的收敛条件。这也说明前面强型界在 $p\downarrow1$ 时至少要有 $(p-1)^{-1}$ 量级的发散。
\end{remark}
% END SUPPLEMENT: id=v14-ch06-ex43-solution
% SOURCE p=223; printed=208; object=ch06-ex44; number=44; region=middle-lower
\begin{exercise}\label{exercise:ch06-fractional-strong-weak}
设 $I_\alpha$ 为\cref{exercise:ch02-fractional-integral}中的分数积分算子。若 $0<\alpha<1$、$1<p<\alpha^{-1}$、$r^{-1}=p^{-1}-\alpha$，则相对于 $(0,\infty)$ 上的 Lebesgue 测度，$I_\alpha$ 为弱 $(1,(1-\alpha)^{-1})$ 型与强 $(p,r)$ 型。
\end{exercise}
% SUPPLEMENT: id=v14-ch06-ex44-solution; parent=ch06-ex44; kind=solution+motivation+hypotheses; origin=independent-derivation
\begin{solution}
\textbf{分数积分核的双侧弱界（重构新增）。}写
\begin{align*}
I_\alpha f(x)=\int_0^\infty K(x,y)f(y)\,\mathrm dy,\qquad
K(x,y)=\frac{(x-y)^{\alpha-1}}{\Gamma(\alpha)}\chi_{\{0<y<x\}}.
\end{align*}
在对角线上将 $K$ 定义为零。令 $q=(1-\alpha)^{-1}>1$。对阈值 $a>0$，令 $L=(\Gamma(\alpha)a)^{-q}$。固定 $x$，$K(x,y)>a$ 等价于 $0<y<x$ 且 $0<x-y<L$，该集合长度为 $\min(x,L)\leqslant L$。固定 $y$，$K(x,y)>a$ 等价于 $y<x<y+L$，其长度正为 $L$。因此两种截面的分布函数均不超过 $(\Gamma(\alpha)a)^{-q}$，所以弱 $L^q$ 准范数均不超过 $C=1/\Gamma(\alpha)$。

两个正半轴均为 $\sigma$ 有限测度空间，核联合可测，故\cref{theorem:ch06-weak-kernel}适用。因 $q'=1/\alpha$，得到弱 $(1,q)$ 型，以及对 $1<p<1/\alpha$ 的强 $(p,r)$ 型，其中
\begin{align*}
\frac1r=\frac1p+\frac1q-1=\frac1p-\alpha.
\end{align*}
同一定理也给出在这些输入指数下定义积分的几乎处处绝对收敛。
\end{solution}
\begin{remark}
\textbf{动机与逐条件说明（重构新增）。}分数积分具有 $(x-y)^{\alpha-1}$ 的局部奇性，其精确水平集长度使弱核定理可以直接使用。$\alpha>0$ 使奇性局部可积，并使 $\Gamma(\alpha)$ 为有限正数；$\alpha<1$ 使 $q>1$ 有限且核随距离递减。$p<1/\alpha$ 保证 $r$ 为有限正数，$p>1$ 则使弱端点插值得到强型。这些条件还给 $r>p$，说明积分改善可积指数。原题没有声称临界 $p=1/\alpha$ 的无穷输出结论，不应从有限指数公式直接推广过去。
\end{remark}
% END SUPPLEMENT: id=v14-ch06-ex44-solution
% SOURCE p=223; printed=208; object=ch06-ex45; number=45; region=lower
\begin{exercise}\label{exercise:ch06-riesz-kernel}
若 $0<\alpha<n$，在 $\mathbb R^n$ 上的函数上定义算子
\begin{align*}
T_\alpha f(x)=\int_{\mathbb R^n}|x-y|^{-\alpha}f(y)\,\mathrm dy.
\end{align*}
原题声称：$T_\alpha$ 相对于 $\mathbb R^n$ 上的 Lebesgue 测度为弱 $(1,(n-\alpha)^{-1})$ 型及强 $(p,r)$ 型，其中 $1<p<n\alpha^{-1}$、$r^{-1}=p^{-1}-\alpha n^{-1}$。（$n=3,\alpha=1$ 的情形在物理中尤其重要：若 $f$ 表示质量或电荷分布的密度，则 $-(4\pi)^{-1}T_1f$ 表示由它产生的引力势或静电势。）
\end{exercise}
% SUPPLEMENT: id=v14-ch06-ex45-solution; parent=ch06-ex45; kind=source-correction+counterexample+solution+motivation+hypotheses; origin=independent-derivation
\begin{remark}
\textbf{题面指数不相容，不能证明原字面结论（重构新增）。}这里已按原页保留核 $|x-y|^{-\alpha}$ 和所列全部指数。这个核的弱指数应为 $n/\alpha$，相应的可积性增益为 $(n-\alpha)/n$；原题把核的奇性指数与分数积分的阶数混用了，而且所列弱指数还缺少维数因子。下面先给出保留原核时的正确结论及证明，再给反例，说明不能把此处当作仅需补写步骤的正确命题。
\end{remark}
\begin{solution}
\textbf{保留原核的正确估计（重构新增）。}令 $q=n/\alpha$、$\omega_n=m(B(0,1))$。对任何固定 $x$ 及 $a>0$，
\begin{align*}
m\{y:|x-y|^{-\alpha}>a\}
=\omega_n a^{-n/\alpha}=\omega_n a^{-q}.
\end{align*}
单个对角点如何赋值不影响测度；固定 $y$ 的列截面同样满足这个公式。因此两侧弱 $L^q$ 准范数恰为 $C=\omega_n^{1/q}=\omega_n^{\alpha/n}$。由 $0<\alpha<n$，$1<q<\infty$，且 $q'=n/(n-\alpha)$。核可测，两空间均 $\sigma$ 有限，由\cref{theorem:ch06-weak-kernel}得到
\begin{align*}
T_\alpha &:L^1\longrightarrow \text{弱 }L^{n/\alpha},\\
T_\alpha &:L^p\longrightarrow L^r,\qquad
1<p<\frac n{n-\alpha},\qquad
\frac1r=\frac1p-\frac{n-\alpha}{n}.
\end{align*}
积分在相应范围几乎处处绝对收敛，算子界只依赖 $n,\alpha,p$。

如果改以 $\beta=n-\alpha$ 表示分数积分的阶数，则同一个核写成 $|x-y|^{\beta-n}$，上述结论变成弱 $(1,n/(n-\beta))$ 型，以及 $1<p<n/\beta$、$1/r=1/p-\beta/n$ 的强型。这说明原题的指数形式如何与核的另一种参数化相对应，但此处没有把原题中的核静默改掉。

\textbf{原强型声称的反例。}取原题中特别提到的 $n=3,\alpha=1$，并取 $p=2$。按原题的指数将得到 $r=6$。令
\begin{align*}
f(y)=|y|^{-7/4}\chi_{\{|y|>1\}}.
\end{align*}
由极坐标，$\int|f|^2=4\pi\int_1^\infty\rho^{2-7/2}\,\mathrm d\rho<\infty$。但对任意固定 $x$，当 $|y|>\max(1,2|x|)$ 时，$|x-y|\leqslant3|y|/2$，因而
\begin{align*}
T_1f(x)\geqslant\frac23\int_{|y|>\max(1,2|x|)}|y|^{-1-7/4}\,\mathrm dy
=\frac{8\pi}3\int_{\max(1,2|x|)}^\infty\rho^{-3/4}\,\mathrm d\rho=\infty.
\end{align*}
所以它甚至处处不收敛，更不可能属于 $L^6$。

\textbf{原弱型声称的反例。}仍取 $n=3,\alpha=1$，输入 $f=\chi_{B(0,1)}\in L^1$。当 $|x|\geqslant2$、$|y|<1$ 时，$|x-y|\leqslant3|x|/2$，故 $T_1f(x)\geqslant c/|x|$，其中 $c=2\omega_3/3>0$。因此对足够小的 $a>0$，
\begin{align*}
\lambda_{T_1f}(a)\geqslant
m\{x:2\leqslant|x|<c/a\}
=\omega_3\bigl((c/a)^3-8\bigr).
\end{align*}
原题所写输出指数为 $(n-\alpha)^{-1}=1/2$；即使采用前节对所有正指数定义的弱空间，也有 $a^{1/2}\lambda_{T_1f}(a)\to\infty$。所以这个弱型结论同样不成立。正确的弱输出指数在此为三。
\end{solution}
\begin{remark}
\textbf{动机与逐条件说明（重构新增）。}距离幂核的水平集是球，维数因子 $n$ 从球体积的幂直接出现。$0<\alpha$ 使核随距离衰减，$\alpha<n$ 使局部奇性可积并给出 $q>1$；正确的 $p<n/(n-\alpha)$ 排除尾部输入使积分发散的情形。将奇性指数与积分阶数分开，是避免本题参数混淆的关键。在 $n=3,\alpha=1$ 时，它是通常阶数为二的势核；原题的物理注释保留其原有归一化，但不能据此推出错误的算子指数。这里的“完成解答”包括反驳错误陈述、给出正确范围和完整证明，而不是给假命题伪造证明。
\end{remark}
% END SUPPLEMENT: id=v14-ch06-ex45-solution
\end{document}
